% !TEX root = ../main.tex
% ---------------------------------------------------------------------
% Apêndice C — as rotinas citadas nominalmente no texto.
%
% MANTIDO À MÃO (ver a nota equivalente no 08_decisoes.tex). Lista apenas as
% rotinas que o corpo cita pelo número, e não as 58 do repositório: o que este
% apêndice resolve é o leitor que encontra "#49" na fonte de uma figura e não
% tem como saber o que é. O índice completo, com as dependências entre
% rotinas, vive no repositório.
%
% Ao acrescentar uma citação de rotina no corpo, acrescentar a linha aqui. O
% verificar.py confere a correspondência nos dois sentidos.
% ---------------------------------------------------------------------
\chapter{Rotinas citadas no texto}
\label{ap:rotinas}

As fontes das figuras e algumas passagens do texto identificam a rotina que
produziu o número pelo seu identificador numérico, na forma
``Pipeline~\#32''. O identificador é permanente e não é reatribuído quando a
ordem lógica do conjunto muda, pelas razões dadas na
Seção~\ref{sec:met-repro}, o que o torna uma referência estável sem torná-la
informativa por si. O Quadro~\ref{quadro:rotinas} diz o que cada rotina
responde e em que arquivo ela está, e cobre as onze que o texto cita
nominalmente, e não as 58 do conjunto, cujo índice completo, com as
dependências, os arquivos de saída e as limitações conhecidas de cada uma,
acompanha o repositório e a visualização interativa (\url{\urlviz}).

Dois identificadores do quadro trazem sufixo de letra, e a convenção merece
registro porque ela é consequência direta da regra de não renumerar. Quando
uma rotina precisa ser inserida no lugar lógico de outra, seja por refazer o
que ela fazia sobre material melhor, seja por estender o que ela deixou em
aberto, a rotina nova recebe o número da antiga acrescido de uma letra, e é
por isso que a matriz de transição recontada sobre o censo é o \#12B e a
decomposição regional da idade é o \#28C. O sufixo diz que existe uma rotina
anterior com aquele número, ainda no repositório e ainda documentada, cujo
resultado a rotina sufixada substituiu ou completou. Apagar a anterior
apagaria a data da correção, que é justamente o que o registro serve para
guardar.

\begin{quadro}[!ht]
\centering
\caption[Rotinas citadas no texto]{Rotinas citadas nominalmente no texto, com a pergunta que cada uma responde.}
\label{quadro:rotinas}
\footnotesize
\begin{tabular}{p{1.3cm}p{6.1cm}p{5.8cm}}
\toprule
\textbf{Rotina} & \textbf{Pergunta que responde} & \textbf{Arquivo em \texttt{scripts/}} \\
\midrule
\#10 & Como era a cobertura do estado em cada um dos quarenta anos? Baixa o
raster do MapBiomas e gera os mapas anuais em resolução nativa
(Figura~\ref{fig:cobertura}) &
\texttt{gerar\_mapas\_lulc\_gee\_40anos.py} \\
\addlinespace
\#12B & Para onde foi cada hectare entre 1985 e 2024? Reconta a matriz de
transição sobre o censo local de pixels, com o Mosaico de Usos como sétimo
grupo (Figura~\ref{fig:sankey}) &
\texttt{transicoes\_cubo.py} \\
\addlinespace
\#25 & Que unidade espacial permite comparar 1985 com 2024? Constrói as 166
AMCs de território constante e o painel que roda sobre elas &
\texttt{construir\_amc\_goias.py} \\
\addlinespace
\#28C & A bimodalidade da idade da pastagem é efeito da região? Decompõe a
distribuição por mesorregião e por AMC sobre o censo
(Figura~\ref{fig:idade}) &
\texttt{bimodalidade\_regional.py} \\
\addlinespace
\#29 & Onde caem as fronteiras entre os períodos? Triangula três métodos de
detecção de quebra e fixa os três atos &
\texttt{periodizacao\_multivariada.py}
\newline \texttt{periodizacao\_stars.py}
\newline \texttt{periodizacao\_transicoes.py} \\
\addlinespace
\#32 & O centro de massa de cada classe se desloca, e em que direção?
Calcula centro médio e centro mediano ano a ano
(Figura~\ref{fig:centromassa}) &
\texttt{centro\_massa.py} \\
\addlinespace
\#38 & O choque cambial incide de forma desigual sobre unidades com exposição
diferente? Estima o desenho de interação no painel de AMCs &
\texttt{drive\_comum\_amc.py} \\
\addlinespace
\#39 & A desaceleração da conversão no Sul vem da demanda ou do esgotamento
da oferta? Decompõe a queda e mede o estoque exposto
(Figura~\ref{fig:fronteira}) &
\texttt{fronteira\_fechando.py} \\
\addlinespace
\#49 & Os canais estimados no painel sobrevivem à dependência espacial?
Estima as especificações com defasagem e com erro espacial
(Figura~\ref{fig:espacial}) &
\texttt{painel\_espacial\_dinamico.py} \\
\addlinespace
\#52 & A exposição do choque pode ser medida por uma característica física,
exógena ao uso da terra? Constrói a camada de aptidão e a substitui no
desenho do \#38 &
\texttt{aptidao\_edafo\_exposicao.py}
\newline \texttt{aptidao\_edafo\_drive38.py} \\
\addlinespace
\#56 & O gradiente estimado é da aptidão ou do eixo Sul--Norte? Põe as
exposições rivais na mesma regressão
(Figura~\ref{fig:horserace}) &
\texttt{drive\_horse\_race\_latitude.py} \\
\bottomrule
\end{tabular}
\fontefig{elaboração própria.}
\end{quadro}
